% !TEX root = ../main.tex
% NB til orchestrator: Denne filen bruker TikZ-figurer i theory-filen optimal-decentralization-model.tex.
% Preamble må utvides med \usepackage{tikz}. Hvis tikz ikke legges til, vil bygget feile på de to TikZ-blokkene i den teorifilen.
% Q06 selv bruker kun \includegraphics og er ikke avhengig av tikz.
\section{Q06: Centralisering vs.\ decentralisering -- trade-offs og optimal grad}\label{q:06}

\begin{quote}
\emph{Redegør for fordele og ulemper ved samt trade-offs mellem centralisering og decentralisering. Præsentér herefter kriterier for fastlæggelse af optimal decentraliseringsgrad og redegør specifikt grafisk for modellen til illustration heraf.}
\end{quote}

\subsection*{Svar (kort)}
\thref{centralization-vs-decentralization}{Sentralisering vs.\ desentralisering} er et \emph{kontinuum} -- en glidende skala, ikke et binært valg. Desentralisering utnytter \thref{specific-vs-general-knowledge}{spesifik lokal viden} (kunnskap som er kostbar å overføre oppover, jf.\ Hayek; Jensen \& Meckling), men introduserer \thref{agency-costs}{agentomkostninger} (kostnaden ved at lokal leder kan handle i egen interesse) og koordineringskostnader. Optimal grad $D^{*}$ ligger der marginal fordel av lokal viden møter marginal kostnad av agent- og koordineringsproblemer. BSZ-modellen (kap.\ 12) gir nettogevinst $\mathrm{NB}(D) = BD - AD - CD^{2}$, dvs.\ desentraliseringsgraden $D$ ganger nytte minus kostnader; jo mer du desentraliserer, jo mer lokal viden ($BD$) -- men også jo mer agentkostnad ($AD$) og en koordineringskostnad som vokser kvadratisk ($CD^{2}$). Optimum er $D^{*} = (B-A)/(2C)$. Tolkning: jo høyere nytten $B$ av lokal viden, jo høyere optimal desentralisering; jo høyere agent- og koordineringskostnad ($A$ og $C$), jo mindre desentralisering. Hver gang beslutningsrettighetene (\SOne) flyttes, må prestasjonsmåling (\STwo) og insentiver (\SThree) justeres med, ellers faller arkitekturen. [SLIDE/BOK]

\subsection*{Fordele og ulemper (kort)}

\begin{center}
\renewcommand{\arraystretch}{1.15}
\begin{tabular}{@{}p{2.5cm}p{6.0cm}p{6.0cm}@{}}
\toprule
& \textbf{Sentralisering} & \textbf{Desentralisering} \\
\midrule
\textbf{Fordele} & Koordinering, stordrift, konsistens, lave agentkostnader & Lokal viden, hurtighet, motivasjon/ledertrening, avlastning av toppledelse \\
\textbf{Ulemper} & Tap av lokal viden, treghet, toppledelsesoverload, demotivasjon nede & Agentkostnader, koordineringssvikt, tap av stordrift, inkonsistente beslutninger \\
\bottomrule
\end{tabular}
\end{center}
\emph{Trade-off i én linje:} Desentralisering gir bedre informasjon men dårligere insentivlinje; sentralisering gir bedre insentivlinje men dårligere informasjon. [SLIDE]

\subsection*{Relevante teorier (kort)}

\paragraph{\thref{centralization-vs-decentralization}{Centralisering vs.\ desentralisering}\quad\SOne}
Beslutningsrett kan plasseres høyt (CEO/konsern) eller delegeres nedover. Valget er det første benet i arkitekturen og det drives av om videnen som trengs til beslutningen er spesifik eller generell. I praksis blandes sentralisering og desentralisering på ulike beslutningsklasser -- f.eks.\ kapitalallokering sentralt, daglig drift lokalt. [SLIDE/BOK]

\paragraph{\thref{optimal-decentralization-model}{Optimal desentraliseringsgrad}\quad\SOne}
BSZ-modellen formaliserer trade-off-en: $\mathrm{NB}(D) = B \cdot D - A \cdot D - C \cdot D^{2}$ med optimum $D^{*} = (B-A)/(2C)$, der $B$ er marginal nytte av lokal viden (gevinst per ekstra desentraliseringssteg), $A$ er marginal agentkostnad og $C$ er konveks koordineringskostnad. Setter man derivatet lik null får man $B - A - 2CD = 0$, og løser for $D^{*}$. \emph{Komparativ statikk} -- altså hvordan optimum endres med parametrene: $\partial D^{*}/\partial B > 0$ (mer verdifull lokal viden $\Rightarrow$ mer desentralisering), $\partial D^{*}/\partial A < 0$ (større agentkostnad $\Rightarrow$ mindre desentralisering), $\partial D^{*}/\partial C < 0$ (større koordineringskostnad $\Rightarrow$ mindre desentralisering). [SLIDE]

\paragraph{\thref{specific-vs-general-knowledge}{Spesifik vs.\ generell viden}\quad\SOne}
Hayek (1945) og Jensen \& Meckling (1992): viden er distribuert. Spesifik viden er kostbar å overføre $\Rightarrow$ flytt \emph{beslutningen} til videnen (desentraliser). Generell viden er billig $\Rightarrow$ flytt \emph{videnen} til beslutningen (sentraliser). Dette er driveren for $B$ i modellen -- og forklarer hvorfor IT-revolusjonen kan re-sentralisere når tidligere spesifik viden kodifiseres. [SLIDE/BOK]

\paragraph{\thref{agency-costs}{Agentomkostninger}\quad\SOne/\STwo/\SThree}
Når beslutningsretten flyttes nedover, mister prinsipalen direkte kontroll. Jensen \& Meckling dekomponerer agentkostnaden: $\text{AC} = \text{monitoring} + \text{bonding} + \text{residual loss}$ -- altså summen av prinsipalens overvåkningsutgifter (monitoring), agentens egne forpliktelsessignaler som covenants og opsjoner (bonding), og det effektivitetstapet som likevel gjenstår (residual loss). Modellens $A$-parameter er nettopp marginal agentkostnad ved økt desentralisering. Den minimeres ved samtidig å designe \STwo\ (hva som måles) og \SThree\ (hvordan agenten belønnes) som erstatter den tapte kontrollen. Endrer man bare \SOne\ uten å justere S2/S3, eksploderer $A$ og $D^{*}$ trekkes mot null. [BOK]

\subsection*{Illustrasjoner fra forelesning}

\begin{figure}[H]
  \centering
  \includegraphics[width=0.6\linewidth]{BSZ_kap_12_2026/by_slide/slide_004/BSZ_kap_12_2026_slide_004_asset_01_image1.jpeg}
  \caption{Sentralisering vs.\ desentralisering -- trade-off mellom lokal viden og kontroll. (BSZ kap.\ 12, slide 4) [SLIDE]}
\end{figure}

\begin{figure}[H]
  \centering
  
\includegraphics[width=0.6\linewidth]{BSZ_kap_12_2026/by_slide/slide_018/BSZ_kap_12_2026_slide_018_asset_01_image6.jpeg}
  \caption{Spesifik viden og kommunikasjon -- hvorfor desentralisering kreves når viden er kostbar å overføre. (BSZ kap.\ 12, slide 18) [SLIDE]}
\end{figure}

\subsection*{Grafisk modell (obligatorisk)}

\paragraph{Fra totalnivå til marginalnivå -- utledning av $D^{*}$.}
Modellen opererer på to nivåer. Først definerer vi \emph{totalkurvene}:
\[
  \mathrm{TB}(D) = B \cdot D \qquad \text{(total fordel -- lineær i $D$)}
\]
\[
  \mathrm{TC}(D) = A \cdot D + C \cdot D^{2} \qquad \text{(total kostnad -- konveks i $D$)}
\]
Nettofordel er differansen mellom total fordel og total kostnad:
\[
  \mathrm{NB}(D) \;=\; \mathrm{TB}(D) - \mathrm{TC}(D) \;=\; BD - AD - CD^{2}
\]
For å finne den $D$ som \emph{maksimerer} NB, deriverer vi NB med hensyn på $D$ og setter lik null:
\[
  \frac{d\,\mathrm{NB}}{dD}
  \;=\; \underbrace{\frac{d(\mathrm{TB})}{dD}}_{=\;\mathrm{MB}\,=\,B}
  \;-\; \underbrace{\frac{d(\mathrm{TC})}{dD}}_{=\;\mathrm{MC}\,=\,A+2CD}
  \;=\; B - A - 2CD \;=\; 0
\]
\[
  \boxed{\;D^{*} = \frac{B-A}{2C}\;}
\]
Altså: å derivere $\mathrm{NB}(D)$ og løse for $D$ gir \emph{direkte} $D^{*}$. Og fordi derivatet av TB er MB og derivatet av TC er MC, er betingelsen $d\mathrm{NB}/dD = 0$ \emph{identisk} med betingelsen $\mathrm{MB} = \mathrm{MC}$. Det er to fremstillinger av \emph{samme} optimum.

\paragraph{Geometrisk kobling mellom de to panelene.}
Figuren nedenfor viser begge nivåer:
\begin{itemize}[nosep]
  \item \textbf{Øvre panel (TC/TB):} NB er den \emph{vertikale avstanden} mellom TB og TC. Denne avstanden er størst ved $D^{*}$.
  \item \textbf{Nedre panel (MC/MB):} $D^{*}$ er skjæringspunktet mellom MB og MC.
\end{itemize}
De to panelene er ekvivalente: der totalkurvene har størst avstand, er også der marginalkurvene krysser -- fordi størst avstand betyr at \emph{stigningsratene} (derivatene) er like.

\begin{figure}[H]
\centering
\begin{tikzpicture}[scale=1.0]
  % === ØVRE PANEL: TC og TB (totalkurver) ===
  \begin{scope}[yshift=7.5cm]
    \draw[->, thick] (0,0) -- (7.5,0) node[right] {$D$};
    \draw[->, thick] (0,0) -- (0,7) node[above] {kr (totalt)};
    % TB = BD = 4D, skalert /4 => D
    \draw[thick, blue!70!black, domain=0:6.8, samples=80]
      plot (\x, {\x}) node[right] {$\mathrm{TB}=BD$};
    % TC = AD + CD^2 = D + 0.5D^2, skalert /4 => 0.25D + 0.125D^2
    \draw[thick, red!70!black, domain=0:6.5, samples=80]
      plot (\x, {0.25*\x + 0.125*\x*\x}) node[right] {$\mathrm{TC}=AD+CD^{2}$};
    % D* striplet linje
    \draw[dashed, gray] (3,0) -- (3,3);
    \node[below] at (3,0) {$D^{*}$};
    % Dobbeltpil for max NB ved D*
    \draw[<->, thick, green!60!black] (3,1.875) -- (3,3);
    \node[right, green!60!black, font=\small] at (3.15,2.44) {$\max\;\mathrm{NB}(D^{*})$};
    % Prikk ved break-even D=6 der TB=TC
    \filldraw[black] (6,6) circle (1.5pt);
    \node[above left, font=\footnotesize] at (6,6) {$\mathrm{TB}=\mathrm{TC}$};
    % Prikker på kurvene ved D*
    \filldraw[blue!70!black] (3,3) circle (1.5pt);
    \filldraw[red!70!black] (3,1.875) circle (1.5pt);
  \end{scope}
  % === Forbindelseslinje mellom panelene ===
  \draw[dashed, gray, thin] (3,4) -- (3,7.5);
  % === NEDRE PANEL: MC og MB (marginalkurver) ===
  \draw[->, thick] (0,0) -- (7.5,0) node[right] {$D$};
  \draw[->, thick] (0,0) -- (0,6) node[above] {kr (marginal)};
  % MB = B = 4, konstant
  \draw[thick, blue!70!black] (0,4) -- (7,4) node[right] {$\mathrm{MB}=B$};
  % MC = A + 2CD = 1 + D
  \draw[thick, red!70!black] (0,1) -- (5,6) node[right] {$\mathrm{MC}=A+2CD$};
  % D* striplet linje
  \draw[dashed, gray] (3,0) -- (3,4);
  \filldraw[black] (3,4) circle (1.5pt);
  \node[below] at (3,0) {$D^{*}$};
  \node[left] at (0,4) {$B$};
  \node[left] at (0,1) {$A$};
  \node[above left, font=\small] at (3,4) {$\mathrm{MB}=\mathrm{MC}$};
\end{tikzpicture}
\caption{Øvre panel: totalkurvene TB (lineær) og TC (konveks) -- størst vertikal avstand ved $D^{*}$. Nedre panel: marginalkurvene MB (horisontal) og MC (stigende) -- krysser ved $D^{*}$. Panelene viser \emph{samme} optimum: maks avstand mellom totalkurvene inntreffer \emph{nøyaktig} der marginalkurvene er like. Talleksempel: $B=4,\; A=1,\; C=0{,}5 \;\Rightarrow\; D^{*}=3$. Skifter $B$ opp $\Rightarrow$ $D^{*}$ til høyre; skifter $A$ eller $C$ opp $\Rightarrow$ $D^{*}$ til venstre. [SLIDE/EGEN]}
\end{figure}

\textbf{Leseveiledning for figuren.} Den horisontale aksen $D$ måler grad av desentralisering fra 0 (full sentralisering) mot høyre (mer desentralisering). I det øvre panelet er TB-linjen rett (stigning $B$) og TC-kurven konveks (stigning $A + 2CD$ som øker i $D$); den grønne pilen markerer $\max\mathrm{NB}$, dvs.\ nettofordelen ved $D^{*}$. I det nedre panelet er MB horisontal på $B$ og MC stiger fra $A$ med stigning $2C$; skjæringspunktet er $D^{*}$. Den stiplede linjen viser at $D^{*}$ i begge paneler er \emph{samme} punkt -- dette er kjernen i modellen: \emph{derivert} $\mathrm{NB}$ lik null $\Leftrightarrow$ $\mathrm{MB} = \mathrm{MC}$. En mer detaljert avledning og en alternativ NB-parabel-figur finnes i teorifilen \thref{optimal-decentralization-model}{optimal decentraliseringsmodell}. [SLIDE]

\subsection*{Real-world eksempler}

\textbf{Eks.\ 1 -- Equinor (E\&P vs.\ konsernfinans).}\quad [EKS]
Letevirksomheten i Equinor er sterkt desentralisert til prosjektledere som investment centers fordi den \thref{specific-vs-general-knowledge}{spesifikke geologiske videnen} ikke kan kodifiseres -- høyt $B$. Samtidig er konsernfinans, kapitalallokering og M\&A sentralisert i Stavanger fordi videnen her er generell (regnskapsstandarder, kapitalkostnad, ratings) og koordineringskostnaden ved fragmentert finans ville vært katastrofalt høy -- høyt $C$. Sammenlignet er $D^{*}_{\text{E\&P}} \gg D^{*}_{\text{finans}}$ innenfor samme konsern, og illustrerer at desentralisering er beslutningsklassespesifikk, ikke firmaspesifikk.

\textbf{Eks.\ 2 -- IKEA vs.\ McDonald's.}\quad [EKS]
IKEA holder produktdesign sentralisert i \"Almhult (designkunnskap er kodifiserbar i tegninger og monteringsanvisninger -- generell viden, lav $B$ ved desentralisering, høy stordrift), men salgsbeslutninger per varehus er desentralisert (lokal kundepreferanse er spesifik). McDonald's tar desentraliseringen lenger gjennom franchising: hele driften er delegert til franchisetakere, mens menyutvalg, brand og forsyningskjede holdes sentralt. Forskjellen mellom de to: McDonald's har bevisst designet \STwo\ (royalty på omsetning) og \SThree\ (residualkrav til franchisetaker) som tillater høyere $D^{*}$ uten at $A$ eksploderer. Sammenligningen viser at $D^{*}$ ikke bare avhenger av bransje, men også av hvor godt arkitekturens to andre ben kan utformes.

\subsection*{Kobling til Mærsk Power-to-X}
PtX-divisjonen er et lærebokeksempel på \emph{differensiert} desentralisering: [CASE]
\begin{itemize}\setlength{\itemsep}{2pt}
\item Teknologivalg (elektrolysørtype, prosessdesign, partnerøkosystem) er bevisst desentralisert fordi \thref{specific-vs-general-knowledge}{spesifik kjemi- og prosessviden} sitter i divisjonen, ikke hos konsernledelsen -- høyt $B$.
\item Kapitalallokering og fundraising holdes \emph{sentralt} fordi A og C der er høye: agentkostnaden ved at en entusiastisk divisjonsleder overinvesterer i tidlig teknologi er stor, og koordineringsbehovet med Mærsks samlede balanse og kredittratings er sterkt.
\item Resultat: ulik $D^{*}$ for ulike beslutningsklasser innenfor samme divisjon -- akkurat som modellen forutsier når $B, A, C$ varierer langs dimensjonene.
\end{itemize}
\textbf{Trade-off:} Hadde konsernet desentralisert også CAPEX-beslutningene uten å styrke \STwo\ (NPV-måling) og \SThree\ (langsiktige insentiver), ville $A$ vokst og $D^{*}_{\text{CAPEX}}$ falt mot null. Dette er presis koblingen til den tre-benede skammelen i Q05.

\subsection*{Fallgruber}
\begin{itemize}\setlength{\itemsep}{2pt}
\item Behandle desentralisering som binær -- tap nyansen at $D$ er kontinuerlig og beslutningsklassespesifikk.
\item Glemme at modellen er marginal -- snakk om \emph{marginal} fordel og kostnad, ikke gjennomsnitt.
\item Bare nevne fordeler/ulemper uten å peke på \emph{driveren} (spesifik vs.\ generell viden) eller \emph{trade-off-en} i marginal-form.
\item Hoppe over den grafiske modellen -- eksamen ber eksplisitt om å ``redegøre specifikt grafisk''.
\item Glemme at $A$ og $C$ er \emph{endogene} -- de avhenger av \STwo\ og \SThree, og kan flyttes ned ved bedre måling og insentivdesign.
\item Forveksle koordineringskostnad ($C$) med agentkostnad ($A$) -- de er to ulike effekter med ulik funksjonsform (kvadratisk vs.\ lineær).
\item Ignorere influence costs og bureaucratic rules som ekstra kostnadsklasse ved sterk sentralisering.
\item Ikke knytte modellen til arkitekturen som helhet -- $D^{*}$ uten S2/S3-justering er meningsløst.
\end{itemize}

\FloatBarrier
\subsection*{Håndskrivbar disposisjon (15-min forberedelse)}
\begin{itemize}\setlength{\itemsep}{2pt}
\item \textbf{Tese:} Optimal grad av desentralisering ligger der marginal fordel av lokal viden = marginal kostnad av agent- og koordineringsproblemer.
\item \textbf{Definisjoner:} Sentralisering (beslutning høyt) vs.\ desentralisering (delegert). Kontinuum, beslutningsklassespesifikt.
\item \textbf{Driver:} \thref{specific-vs-general-knowledge}{Spesifik vs.\ generell viden} (Hayek $\to$ J\&M). Co-location of knowledge and decision rights.
\item \textbf{Fordeler ved desentralisering:} Lokal viden, hurtighet, toppledertid, motivasjon/ledertrening.
\item \textbf{Ulemper ved desentralisering:} \thref{agency-costs}{Agentkostnader}, koordineringskostnader, tap av stordrift.
\item \textbf{Speilbilde:} Sentralisering har omvendte fordeler/ulemper -- viktig å nevne begge.
\item \textbf{Trade-off:} Bedre info vs.\ bedre insentivlinje. Flytter man \SOne, må \STwo\ og \SThree\ med -- ellers eksploderer $A$.
\item \textbf{Modell:} $\mathrm{NB}(D) = BD - AD - CD^{2}$. FOC: $B - A - 2CD = 0 \Rightarrow D^{*} = (B-A)/(2C)$. SOC: $-2C < 0$ (maks).
\item \textbf{Komparativ statikk:} $\partial D^{*}/\partial B > 0$; $\partial D^{*}/\partial A < 0$; $\partial D^{*}/\partial C < 0$.
\item \textbf{Graf (to paneler):} Øvre panel: TB=$BD$ (rett linje) og TC=$AD+CD^{2}$ (konveks kurve); vis at størst avstand er ved $D^{*}$. Nedre panel: MB=$B$ (horisontal) og MC=$A+2CD$ (stigende); skjæringspunkt = $D^{*}$. Poenget: maks avstand oppe = kryss nede $\Leftrightarrow$ derivert NB = 0.
\item \textbf{Talleksempel:} $B=10, A=2, C=4 \Rightarrow D^{*} = 1$. Endre $A$ til $6 \Rightarrow D^{*} = 0.5$.
\item \textbf{Eks.\ 1:} Equinor E\&P (høy $D^{*}$) vs.\ konsernfinans (lav $D^{*}$) i samme firma.
\item \textbf{Eks.\ 2:} IKEA (design sentralt, salg lokalt) vs.\ McDonald's franchise (sterk desentralisering støttet av sterk \STwo/\SThree).
\item \textbf{Case:} Mærsk PtX -- teknologi desentralisert, CAPEX sentralisert; viser at $D^{*}$ er beslutningsklassespesifikk.
\item \textbf{Fallgruber:} Ikke gjør $D$ binær; husk marginal; knytt $A$ til S2/S3-kvalitet; tegn grafen.
\item \textbf{Søyle:} Hovedsakelig \SOne, men kobler direkte til \STwo\ og \SThree\ via parameteren $A$.
\end{itemize}
